\section{DDD y contextos delimitados en un dominio de seguros}

\subsection{Enunciado}
Una empresa de seguros necesita desarrollar un sistema para gestionar pólizas de automóviles. El
sistema debe interactuar con un departamento de siniestros, un sistema de facturación de primas
y una aplicación para que los ajustadores de seguros evalúen daños en el lugar del accidente.
Cada uno de estos subsistemas tiene su propia terminología y preocupaciones (por ejemplo,
«póliza» para un actuario es un contrato con cláusulas, para facturación es una fuente de
ingresos, y para un ajustador es un número de referencia).

Basado en los conceptos de Domain-Driven Design (Evans) y Bounded Context (Richards \& Ford),
explique cómo dividiría este sistema en varios contextos delimitados. ¿Cómo evitaría que el
modelo de «Póliza» se contamine con detalles de siniestros o facturación? ¿Cómo se comunicarían
estos contextos de manera desacoplada? Relacione esta solución con la implementación de una
Arquitectura Orientada al Dominio (como Onion o Clean) y la creación de un Adaptador de
Servicio para la comunicación externa, tal como se describe en la Guía de Arquitectura N-Capas.

\subsection{Respuesta}
% TODO: bounded contexts; evitar contaminación del modelo; comunicación desacoplada;
% relación con Onion/Clean y adaptadores de servicio (guía N-Capas).
\textit{[Redactar con \parencite{richards2020fundamentals,torre2010guia}.]}
